% Independent raw-output obstruction module, 18 November 2026.
% Construction suggested by root; saturation estimate independently checked
% by fresh_reader_approximation and robust_agnostic.
\section{A raw posterior mean can fail under finite-variance Student weights}

Fix $p>2$ and $T>0$. Let the three parameters $v,w,\gamma$ be
independent with density
\[
 q_p(t)=c_p(1+t^2)^{-(p+1)/2},\qquad t\in\mathbb R,
\]
and consider the ordinary biased one-neuron network
$f_{v,w,\gamma}(x)=v\tanh(wx+\gamma)$.
Every coordinate has a finite second moment. The prior, architecture,
and likelihood temperature are fixed independently of the sample size.
For observations $S_n$, use the raw Gaussian likelihood and raw mean
\[
 dQ_{S_n}(\theta)=Z_{S_n}^{-1}
  \exp\left\{-\frac1T\sum_{i=1}^n(f_\theta(X_i)-Y_i)^2\right\}
  d\pi(\theta),\qquad
 \bar f_{S_n}(x)=\int f_\theta(x)\,dQ_{S_n}(\theta).
\]
The physical observation variance is $T/2$. There is no clipping.

\begin{theorem}[A distribution-uniform raw-mean obstruction]
There are constants $c_*,q_*>0$, depending only on $p,T$, and
distributions $P_n$ supported on $[-1,1]\times[-1,1]$, with noiseless
labels, such that for every sufficiently large $n$,
\[
 P_n^{\otimes n}\{R_{P_n}(\bar f_{S_n})\ge c_*n^{1/3}\}\ge q_*.
\]
Yet the fixed width-one comparator $g\equiv0$ satisfies
$R_{P_n}(g)\le n^{-1/3}$. In particular,
$\mathbb E_{P_n^{\otimes n}}R_{P_n}(\bar f_{S_n})\ge q_*c_*n^{1/3}$.
Thus this fixed prior and fixed likelihood do not admit any
distribution-uniform vanishing oracle remainder against the zero
network, either in expectation or at any fixed failure probability
$\delta<q_*$.
\end{theorem}

\begin{proof}
Write
\[
 b=n^{-1/6},\quad c=\frac1{100p(T+1)},\quad
 R=e^{cn^{2/3}},\quad a=b/R.
\]
The distribution assigns probabilities $(1-1/n)/2,(1-1/n)/2,1/n$
to the pairs $(a,b),(-a,-b),(1,0)$, respectively.
Let $E_n$ be the event that the sample contains no input $1$ and that
its counts $n_+,n_-$ at $a,-a$ both belong to $[n/3,2n/3]$.
Its probability tends to $e^{-1}$, so is bounded below by a positive
constant. All bounds below are uniform over samples in $E_n$.

For such a sample define
\[
 u_+=v\tanh(\gamma+aw),\quad
 u_-=v\tanh(\gamma-aw),\quad
 L=\exp\{-[n_+(u_+-b)^2+n_-(u_-+b)^2]/T\}.
\]
The transformation $(v,w,\gamma)\mapsto(-v,-w,-\gamma)$ leaves
the represented function unchanged and preserves the prior. We may
therefore restrict all integrals to $v>0$ and multiply by two.
Constants in the proof may depend on $p,T$ but never on $n$ or the
balanced counts.

Set
\[
 G=\{|u_+-b|\le b/2,\ |u_-+b|\le b/2\},\qquad
 Z_*:=\frac{a^p b^{-p-1}}n=\frac{R^{-p}}{bn}.
\]
On $G^c$, the likelihood is at most
\begin{equation}\label{eq:raw-student-bad-fit}
 L\le \exp\{-nb^2/(12T)\}=\exp\{-n^{2/3}/(12T)\}.
\end{equation}
Since $pc=1/[100(T+1)]<1/(12T)$, this upper bound is
$o(Z_*)$, and also $o(Z_*\log R)$.

We first prove the evidence bound $Z_{S_n}\le C Z_*$.
On $G$ necessarily $v\ge b/2$. For $b/2\le v\le3b$, set
$h_+=\tanh(\gamma+aw)$ and $h_-=-\tanh(\gamma-aw)$.
If the good-fit set of such $v$ is nonempty, then
$h_+,h_-\ge1/6$, whence $w\ge t_0/a$ for
$t_0=\operatorname{arctanh}(1/6)>0$.
For every fixed such $(w,\gamma)$, the exponent as a function of $v$
is a quadratic with coefficient
\[
 n_+h_+^2+n_-h_-^2\ge n/36.
\]
The density $q_p(v)$ is bounded, so integrating in $v$ first, even
over the entire real line, gives at most $C/\sqrt n$.
The Student tail bound then gives
\begin{equation}\label{eq:raw-student-saturation}
 \int_{G\cap\{b/2\le v\le3b\}}L\,d\pi
 \le\frac C{\sqrt n}\int_{t_0/a}^\infty q_p(w)\,dw
 \le \frac{C a^p}{\sqrt n}.
\end{equation}
The bias density integrates to one. This argument includes the
near-saturation region without a singular change of variables.
The ratio of the last bound to $Z_*$ is
$b^{p+1}\sqrt n=n^{(2-p)/6}\le1$ for large $n$.

For $v\ge3b$, use the exact coordinates $(u_+,u_-)$ instead of
$(w,\gamma)$:
\[
 w=\frac{\operatorname{arctanh}(u_+/v)
             -\operatorname{arctanh}(u_-/v)}{2a},\qquad
 \gamma=\frac{\operatorname{arctanh}(u_+/v)
             +\operatorname{arctanh}(u_-/v)}2.
\]
Their Jacobian is
\begin{equation}\label{eq:raw-student-jacobian}
 \left|\frac{\partial(w,\gamma)}{\partial(u_+,u_-)}\right|
 =\frac1{2av^2[1-(u_+/v)^2][1-(u_-/v)^2]}.
\end{equation}
On $G$ we have $|u_\pm|/v\le1/2$ and
$R/(2v)\le w\le2R/v$. Thus the Jacobian is at most $C/(av^2)$.
The global bound $q_p(w)\le C|w|^{-p-1}$ and boundedness of
$q_p(\gamma)$ imply
\[
 q_p(v)q_p(w)q_p(\gamma)
 \left|\frac{\partial(w,\gamma)}{\partial(u_+,u_-)}\right|
 \le C a^p b^{-p-1}q_p(v)v^{p-1}.
\]
The $v$ integral of $q_p(v)v^{p-1}$ over $(0,\infty)$ is finite.
The likelihood integral over the $u$ rectangle is at most
$\pi T/\sqrt{n_+n_-}\le C/n$.
Together with \eqref{eq:raw-student-bad-fit} and
\eqref{eq:raw-student-saturation}, this proves
\begin{equation}\label{eq:raw-student-evidence}
 Z_{S_n}\le C Z_*.
\end{equation}

We next lower-bound the numerator of the raw prediction at $x=1$.
Let $t_n=1/(10\sqrt n)$ and restrict the transformed integral to
\[
 1\le v\le R/100,\quad |u_+-b|\le t_n,\quad
 |u_-+b|\le t_n.
\]
For large $n$ this region lies inside $G$; its likelihood is at least
$e^{-1/(100T)}$. Here $R/(2v)\le w\le2R/v$, $|\gamma|\le1$,
and hence $f(1)=v\tanh(w+\gamma)\ge c_1v$.
The Student density lower bounds give
\[
 q_p(v)\ge c_2v^{-p-1},\quad
 q_p(w)\ge c_2(R/v)^{-p-1},\quad q_p(\gamma)\ge c_2.
\]
The Jacobian in \eqref{eq:raw-student-jacobian} is at least
$1/(2av^2)$. The prediction-weighted density is consequently at
least $c_3 a^p b^{-p-1}/v$ throughout this region. Its $u$ area is
$4t_n^2=1/(25n)$, and its $v$ integral contains
$\int_1^{R/100}dv/v=\log(R/100)$. Its positive contribution to the
numerator is therefore at least $c_4 Z_*\log R$.

There is no negative contribution on $G$. Indeed, for $v>0$ the
opposite signs of $u_+,u_-$ imply $w>0$ and $-aw<\gamma<aw$.
Since $a<1$, it follows that $w+\gamma>w(1-a)>0$.
On $G^c$, $|f(1)|\le|v|$ and
\eqref{eq:raw-student-bad-fit} bounds the absolute contribution by
$\mathbb E_\pi|v|e^{-n^{2/3}/(12T)}=o(Z_*\log R)$.
After symmetry is restored, the full numerator is still at least
$c_5Z_*\log R$. By \eqref{eq:raw-student-evidence},
\[
 \bar f_{S_n}(1)\ge c_6\log R=c_6c n^{2/3}
 \qquad(S_n\in E_n).
\]
The rare input alone contributes at least $(c_6c)^2n^{1/3}$ to test risk.
The zero comparator has risk $(1-1/n)b^2\le n^{-1/3}$, completing
the proof. All posterior quantities exist: the likelihood is bounded
by one, its normalizer is positive, and $|f(x)|\le|v|$ with
$\mathbb E_\pi v^2<\infty$.
\end{proof}

\begin{corollary}[A cheap comparison law still does not give a raw oracle]
Let $\rho_n=\pi(\cdot\mid |v|\le1/n)$. For every data distribution
in the theorem,
\[
 \operatorname{KL}(\rho_n\|\pi)\le\log n+C_p,
 \qquad
 \int R_{P_n}(f)\,d\rho_n(f)
       \le R_{P_n}(0)+n^{-2}.
\]
Consequently no fixed constant $A<\infty$ can make the naive raw-mean
PAC--Bayesian bound
\[
 \mathbb E R_{P_n}(\bar f_{S_n})
 \le\inf_{\rho\ll\pi}
 \left\{\int R_{P_n}(f)\,d\rho(f)
                 +\frac{AT}{n}\operatorname{KL}(\rho\|\pi)\right\}
\]
valid for all bounded data distributions and all $n$ for this prior
and fixed temperature.
\end{corollary}
\begin{proof}
The density $q_p$ is bounded below on $[-1,1]$, so
$\pi(|v|\le1/n)\ge c_p'/n$, giving the entropy bound.
Conditioning on this event preserves the symmetry and independence
of the readout $v$. Thus $\int f(x)d\rho_n=0$ and
$\int f(x)^2d\rho_n\le n^{-2}$ at every input.
Expanding squared loss gives the risk bound. The displayed
PAC--Bayesian right side would then tend to zero, whereas the main
theorem makes its left side diverge. The comparison measure is
deterministic and uses no training data.
\end{proof}

\paragraph{Scope.}
The adverse distribution depends on $n$. This is a failure of a
distribution-uniform oracle theorem, not a claim of inconsistency for
one fixed distribution. The common inputs approach each other
exponentially fast, so the construction imposes no fixed geometric
regularity or uniform anti-concentration constant on the input law.
The teacher is the same zero function for every $n$. The prior is
dense, continuous, independent across all three coordinates, and
unchanged with $n$; the likelihood is the ordinary all-data Gaussian
likelihood with fixed variance. The lower bound is specifically for
the raw posterior mean, rather than for posterior-averaged loss.
